% ----- DOHA 2 -----

\newpage

\subsection*{\deva{द्वितीय दोहा} (Second Doha)}
\addcontentsline{toc}{subsection}{\deva{द्वितीय दोहा} (Second Doha) — \deva{बुद्धिहीन तनु...}}

\begin{minipage}[t]{0.55\textwidth}
\vspace{0pt}

{\large \linespread{1.5}\selectfont
\makebox[\linewidth][l]{\deva{बुद्धिहीन तनु जानिके}}\\
\makebox[\linewidth][c]{\deva{सुमिरौं पवन कुमार।}}\\
\makebox[\linewidth][l]{\deva{बल बुधि बिद्या देहु मोहिं }}\\
\makebox[\linewidth][c]{\deva{हरहु कलेस बिकार॥}}
\par}

\vspace{0.5em}

{\large \linespread{1.5}\selectfont
\makebox[\linewidth][l]{\telu{బుద్ధిహీన తను జానికే}}\\
\makebox[\linewidth][c]{\telu{సుమిరౌఁ పవన కుమార।}}\\
\makebox[\linewidth][l]{\telu{బల బుధి బిద్యా దేహు మోహిఁ}}\\
\makebox[\linewidth][c]{\telu{హరహు కలేస బికార॥}}
\par}

\vspace{0.5em}

{\large \linespread{1.3}\selectfont
\textit{buddhihīna tanu jānike, sumiraũ pavana kumāra |}\\
\textit{bala budhi bidyā dehu mohĩ, harahu kalesa bikāra ||}
\par}
\end{minipage}%
\hfill
\begin{minipage}[t]{0.42\textwidth}
\vspace{0pt}
\begin{tcolorbox}[anchorbox]
\textbf{The Qualities of Success:} We ask for strength (\textit{Bala}), intellect (\textit{Budhi}), and wisdom (\textit{Bidya}). In the *Valmiki Ramayana* (*Sundara Kanda, 1:198*), the celestials praise Hanuman, declaring that he who has four key leadership traits—courage (\textit{Dhriti}), vision (\textit{Drishti}), intellect (\textit{Mati}), and skill (\textit{Dakshya})—like him, will never fail in his undertakings.
\vspace{0.5em}\newline Hanuman's four timeless qualities of courage, vision, intellect, and skill represent the ultimate blueprint for overcoming life's challenges.\newline\null\hfill\href{https://www.valmiki.iitk.ac.in/sloka?field_kanda_tid=5&language=dv&field_sarga_value=1}{\textit{Sundara Kanda, 1:198}}
\end{tcolorbox}
\end{minipage}

\vspace{0.5em}
\subsubsection*{\textbf{Visual Rhythm Map:}}
\begin{table}[h!]
\centering
\renewcommand{\arraystretch}{1.2}
\begin{tabularx}{\textwidth}{|>{\centering\arraybackslash\hsize=1.5000\hsize}X|>{\centering\arraybackslash\hsize=0.5000\hsize}X|>{\centering\arraybackslash\hsize=1.2500\hsize}X|>{\centering\arraybackslash\hsize=1.0000\hsize}X|>{\centering\arraybackslash\hsize=0.7500\hsize}X|>{\centering\arraybackslash\hsize=1.0000\hsize}X|}
\hline
\textbf{\guru{बु}\laghu{द्धि}\guru{ही}\laghu{न}} \par (6) & \textbf{\laghu{त}\laghu{नु}} \par (2) & \textbf{\guru{जा}\laghu{नि}\guru{के}} \par (5) & \textbf{\laghu{सु}\laghu{मि}\guru{रौ}} \par (4) & \textbf{\laghu{प}\laghu{व}\laghu{न}} \par (3) & \textbf{\laghu{कु}\guru{मा}\laghu{र}} \par (4) \\
\hline
\end{tabularx}
\vspace{-1.5em}
\end{table}

\begin{table}[h!]
\centering
\renewcommand{\arraystretch}{1.2}
\begin{tabularx}{\textwidth}{|>{\centering\arraybackslash\hsize=0.6667\hsize}X|>{\centering\arraybackslash\hsize=0.6667\hsize}X|>{\centering\arraybackslash\hsize=1.0000\hsize}X|>{\centering\arraybackslash\hsize=1.0000\hsize}X|>{\centering\arraybackslash\hsize=1.0000\hsize}X|>{\centering\arraybackslash\hsize=1.0000\hsize}X|>{\centering\arraybackslash\hsize=1.3333\hsize}X|>{\centering\arraybackslash\hsize=1.3333\hsize}X|}
\hline
\textbf{\laghu{ब}\laghu{ल}} \par (2) & \textbf{\laghu{बु}\laghu{धि}} \par (2) & \textbf{\laghu{बि}\guru{द्या}} \par (3) & \textbf{\guru{दे}\laghu{हु}} \par (3) & \textbf{\guru{मो}\laghu{हिं}} \par (3) & \textbf{\laghu{ह}\laghu{र}\laghu{हु}} \par (3) & \textbf{\laghu{क}\guru{ले}\laghu{स}} \par (4) & \textbf{\laghu{बि}\guru{का}\laghu{र}} \par (4) \\
\hline
\end{tabularx}
\vspace{-1.5em}
\end{table}

\vspace{0.5em}
\textbf{ \deva{सरल-संस्कृतम्}:}\\
 \deva{हे पवनकुमार! मां बुद्धिहीनं ज्ञात्वा (मत्वा) अहं त्वां स्मरामि।}\\
 \deva{मह्यं बलं, बुद्धिं, विद्यां च यच्छ (देहि), मम सर्वान् क्लेशान् विकारान् (दोषान्) च दूरीकुरु।}\\

Knowing myself to be devoid of intelligence, O Son of the Wind, I remember you.\\
Grant me strength, wisdom, and knowledge, and remove all my afflictions and impurities.


\vspace{1em}
\newpage
\textbf{\deva{प्रतिपदार्थः} (Meaning of each word):}
\begin{table}[h!]
\centering
\renewcommand{\arraystretch}{1.2}
\begin{tabularx}{\textwidth}{@{} l l X X @{}}
\toprule
\textbf{Awadhi} & \textbf{Sanskrit} & \textbf{English} & \textbf{Grammar \& Notes} \\
\midrule
\deva{बुद्धिहीन} / \telu{బుద్ధిహీన} / \textit{buddhihīna}  &  \deva{बुद्धि-रहितं}  &  Devoid of intelligence  &   \\
\deva{तनु} / \telu{తను} / \textit{tanu}  &  \deva{शरीरं / आत्मानम्}  &  Body / Myself  &   \\
\deva{जानिके} \gotchaicon / \telu{జానికే} / \textit{jānike}  &  \deva{ज्ञात्वा}  &  Having known  & Conjunctive Participle (-ike) \\
\deva{सुमिरौं} / \telu{సుమిరౌఁ} / \textit{sumiraũ}  &  \deva{स्मरामि}  &  I remember/meditate  &   \\
\deva{पवन-कुमार} / \telu{పవన-కుమార} / \textit{pavana kumāra}  &  \deva{हे वायुपुत्र}  &  O Son of the Wind  &   \\
\deva{बल} / \telu{బల} / \textit{bala}  &  \deva{बलम्}  &  Strength  &   \\
\deva{बुधि} \gotchaicon / \telu{బుధి} / \textit{budhi}  &  \deva{बुद्धिम्}  &  Wisdom  & Meter tetris (1+1=2) \\
\deva{बिद्या} / \telu{బిద్యా} / \textit{bidyā}  &  \deva{विद्याम्}  &  Knowledge  & \\ 
\deva{देहु} \gotchaicon / \telu{దేహు} / \textit{dehu}  &  \deva{देहि (यच्छ)}  &  Give / Grant  & Imperative -hu ending \\
\deva{मोहिं} \gotchaicon / \telu{మోహిఁ} / \textit{mohĩ}  &  \deva{मह्यम्}  &  To me  & Dative marker (-hiṃ) \\
\deva{हरहु} \gotchaicon / \telu{హరహు} / \textit{harahu}  &  \deva{हर (दूरीकुरु)}  &  Remove / Destroy  & Imperative -hu ending \\
\deva{कलेस} / \telu{కలేస} / \textit{kalesa}  &  \deva{क्लेशान् / दुःखानि}  &  Sufferings  & \\ 
\deva{बिकार} / \telu{బికార} / \textit{bikāra}  &  \deva{विकारान् / दोषान्}  &  Impurities / Flaws  & \\ 
\bottomrule
\end{tabularx}
\end{table}

\begin{tcolorbox}[gotchabox]
\begin{itemize}
    \item \textbf{Verbs ending in -hu (\deva{देहु, हरहु}):} Adding -hu (-ahu) to a root is the Awadhi equivalent of the Sanskrit imperative mood (making a request/command).
    \item \textbf{Dative Suffix -hiṃ (\deva{मोहिं}):} Maps to Sanskrit Mahyam ("to me"). "Dehu mohĩ" = Give to me.
    \item \textbf{Anusvara vs Chandrabindu (\deva{मोहिं} / \telu{మోహిఁ}):} Why does Devanagari use a dot (\deva{ं}) while Telugu uses a half-circle (\telu{ఁ}) for the exact same nasal sound? In Devanagari orthography, when a vowel matra extends \textit{above} the top line (like \deva{ि, ी, े, ो}), there is physically no vertical space to draw the full Chandrabindu (\deva{ँ}). So, purely for typographical convenience, writers use the Anusvara dot (\deva{ं}) even though it is pronounced as a nasalized vowel. However, if the vowel matra goes \textit{below} or beside the line (like the 'u' in \deva{तिहुँ}), the space above is empty, so the full Chandrabindu \deva{ँ} is correctly written! Telugu script has no top-line constraint, so it always correctly uses the Ara-sunna (\telu{ఁ}).
    \item \textbf{Meter Tetris (\deva{बुद्धि} vs \deva{बुधि}):} Line 1 heavily weights \deva{बुद्धि} (bu + conjunct ddhi) to make it 2+1+2+1=6 beats. Line 2 uses light Awadhi \deva{बुधि} (1+1=2 beats) to save space!
    \item \textbf{The 'e' vowel weight (\deva{जानिके}):} The sound 'e' is historically a compound vowel (a + i = e) and physically requires twice as much breath to produce, making it an automatic 2-beat Guru.
\end{itemize}
\end{tcolorbox}

\newpage
